\documentclass{article}
\usepackage{neurips_2026}
\usepackage[utf8]{inputenc}
\usepackage[T1]{fontenc}
\usepackage{hyperref}
\usepackage{url}
\usepackage{booktabs}
\usepackage{amsfonts,amsmath,amssymb,amsthm}
\usepackage{nicefrac}
\usepackage{microtype}
\usepackage{xcolor}
\usepackage{graphicx}

\newtheorem{proposition}{Proposition}
\newtheorem{corollary}{Corollary}
\newcommand{\E}{\mathbb{E}}
\newcommand{\Var}{\mathrm{Var}}
\newcommand{\Cov}{\mathrm{Cov}}
\newcommand{\ppipp}{PPI++}

\title{How Much Shorter Can a Prediction-Powered Confidence Interval Be?\\ A Headroom Analysis with Fixed Labels, Unlabeled Data and Predictor}

\author{Anonymous Author(s)}

\begin{document}
\maketitle

\begin{abstract}
Prediction-powered inference (PPI) and its power-tuned variant PPI++ combine a small labeled sample, a large unlabeled sample and a fixed pretrained predictor into valid confidence intervals. We ask a narrow question: \emph{with all three inputs held fixed and nominal coverage required, how much shorter can the interval for a population mean still get, and when?}
We give an exact two-sample variance formula for the augmented estimator $\hat\theta_g$ (any fixed $g$ of what is observed on every unit), show that the optimally power-tuned PPI++ interval is exactly the best \emph{affine-in-the-prediction} member of this family, and derive the remaining headroom in closed form: it equals the variance of the part of $\E[Y\mid\text{observables}]$ that is not affine in the prediction, scaled by $N/(n(n+N))$. For regular estimators this is the whole headroom (semiparametric efficiency bound), so the analysis also says when nothing more can be gained.
We check this on nine public datasets (ecology, astronomy, census, proteomics, elections, genomics; 1000 replications per cell) and eight synthetic designs with known truth, comparing PPI, PPI++, a classical interval and cross-fitted augmented estimators with affine, isotonic, covariate-linear and boosted-tree calibration. On all nine real datasets the estimated headroom is at most $2.4\%$ of the variance and the best cross-fitted interval is within $[0.967,1.011]$ of the PPI++ width at $n=200$ labels (within $[0.996,1.002]$ at $n=1000$); the only visible gain ($\approx3$--$7\%$) comes from a predictor whose slope exceeds one. On synthetic designs where the predictor is non-linearly calibrated or ignores informative covariates, widths fall to $0.49$--$0.61$ of PPI++ and match the formula. The larger practical issue is validity: Wald intervals of all methods under-cover at small $n$ (e.g.\ PPI++ $76\%$ at $n=25$ on census income), a studentized bootstrap repairs only part of it, and under covariate shift between labeled and unlabeled data the power-tuned intervals fail completely ($0\%$ coverage) while $\lambda=1$ PPI does not.
\end{abstract}

\section{Introduction}
In many sciences gold-standard labels are scarce, while features and a pretrained predictor are cheap. Prediction-powered inference \citep{angelopoulos2023ppi} uses the predictor on a large unlabeled sample and corrects its bias with a small labeled sample, yielding intervals that are valid regardless of predictor quality and shorter than labeled-only intervals when the predictor is good. PPI++ \citep{angelopoulos2023ppipp} adds a data-driven weight $\lambda$ on the predictions so the interval never gets much worse than the classical one. Subsequent work changes the inputs: cross-fitting when the predictor is trained on the labels \citep{zrnic2023cross}, choosing which points to label \citep{zrnic2024active}, stratification \citep{fisch2024stratified}.

We study the complementary question in which \emph{none of the inputs may change}: the labeled set, the unlabeled set and the predictor are given, and only the way they are combined is ours to choose. Is the PPI++ interval already as short as a valid interval can be? Our contributions are:
\begin{enumerate}
\item \textbf{A closed-form headroom (Section~\ref{sec:theory}).} An exact two-sample variance for the augmented estimator, the identification of optimally tuned PPI++ with its best affine member, and the exact gap to the efficient estimator: $V_{\rm PPI++}-V_{\rm eff}=\frac{N}{n(n+N)}\Var(\mu(Z)-b^\star F)$. The gap is zero when the calibration curve is affine and no covariate carries information the predictor lacks, and it vanishes as $N/n\to0$. Most ingredients are classical (missing-data and semi-supervised mean estimation \citep{robins1994estimation,zhang2019semisup}); the point is the exact comparison with PPI++ and the use of it as a diagnostic.
\item \textbf{A measurement of the headroom on real data (Section~\ref{sec:real}).} For all nine datasets in the PPI reference package the headroom for the mean is $\le2.4\%$ of the variance, and cross-fitted estimators do not shorten intervals appreciably; we report where the small gains are.
\item \textbf{Conditions under which large gains exist and are realised (Section~\ref{sec:synth}).} Synthetic designs with a non-linear calibration curve or an informative covariate yield intervals $40$--$50\%$ shorter, matching the formula once $n$ is moderate; the price at small $n$ is quantified.
\item \textbf{Validity caveats (Sections~\ref{sec:valid},\ref{sec:shift}).} Wald intervals of every method are anti-conservative at small $n$ on real data; shorter is not better unless coverage is checked. Under labeled/unlabeled shift, the gains that come from fitting calibration on labeled data turn into complete failure.
\end{enumerate}
The honest summary is conditional: for a population mean with a reasonable predictor, PPI++ leaves little to gain; there is room only when the predictor is miscalibrated beyond an affine map or leaves covariate information unused.

\section{Related Work}
\textbf{PPI family.} PPI \citep{angelopoulos2023ppi} and PPI++ \citep{angelopoulos2023ppipp} give asymptotically valid intervals for means, quantiles and M-estimands; PPI++ tunes a scalar (or matrix) weight on the predictor. Cross-PPI \citep{zrnic2023cross} handles predictors trained on the labeled data by cross-fitting; active inference \citep{zrnic2024active} changes which units are labeled; stratified PPI \citep{fisch2024stratified} applies PPI within strata of the predictor score, which corresponds to a piecewise-constant $g$ in our notation. These change inputs or use particular calibration classes; we instead ask what the best use of fixed inputs could achieve and measure it.\\
\textbf{Semiparametric efficiency, semi-supervised and post-prediction inference.} Estimating a mean when outcomes are missing completely at random is a classical semiparametric problem whose efficient estimators are augmented inverse-probability-weighted \citep{robins1994estimation,bickel1993efficient,tsiatis2006semiparametric}, with semi-supervised mean estimation developed in \citet{zhang2019semisup}. Machine-learned nuisance functions are handled by cross-fitting \citep{chernozhukov2018dml}. POP-Inf \citep{miao2023popinf} proposes an assumption-lean, data-adaptive post-prediction estimator that is at least as efficient as the classical and PPI estimators. Our estimators are standard members of this family; our contribution is the exact PPI++ comparison, the headroom decomposition and its empirical evaluation (including finite-sample validity and shift), not a new estimator.

\section{Setting and a Headroom Theory}\label{sec:theory}
\textbf{Setting.} A predictor $f$ is fixed. Every unit has observables $Z=(F,W)$ with $F=f(\text{features})$ the prediction and $W$ optional covariates (empty in some datasets). We observe $n$ labeled units $(Z_i,Y_i)\overset{iid}{\sim}P$ and $N$ unlabeled units $Z'_j\overset{iid}{\sim}P_Z$, independent. The target is $\theta=\E[Y]$; $\mu(Z)=\E[Y\mid Z]$. We choose the mean because it is the canonical PPI target, has a closed-form efficiency theory, and every dataset has a natural mean (Section~\ref{sec:setup}). The assumption that labeled and unlabeled units are exchangeable (A1: the i.i.d.\ sampling just described) is what the validity of all methods rests on; Section~\ref{sec:shift} shows what happens without it.

For a fixed function $g$ with $\E g(Z)^2<\infty$ define the augmented estimator
\begin{equation}\label{eq:aug}
\hat\theta_g=\frac1{n+N}\Big(\sum_{i\le n}g(Z_i)+\sum_{j\le N}g(Z_j')\Big)+\frac1n\sum_{i\le n}\big(Y_i-g(Z_i)\big).
\end{equation}
\begin{proposition}[exact variance]\label{prop:var}
$\hat\theta_g$ is unbiased and $\;\Var(\hat\theta_g)=\dfrac{\Var(Y)}{n+N}+\dfrac{N}{n(n+N)}\Var\big(Y-g(Z)\big).$ The minimum over $g$ is attained at $g=\mu$ (up to constants): $V_{\rm eff}=\frac{\Var(Y)}{n+N}+\frac{N}{n(n+N)}\E\Var(Y\mid Z)$.
\end{proposition}
\begin{proposition}[PPI++ is the best affine member]\label{prop:ppi}
For $\lambda\in\mathbb R$ let $\hat\theta^{\rm PPI}_\lambda=\frac1n\sum_i(Y_i-\lambda F_i)+\frac\lambda N\sum_jF'_j$, with variance $\Var(Y-\lambda F)/n+\lambda^2\Var(F)/N$. Its minimum over $\lambda$ (attained at $\lambda^\star=\Cov(Y,F)/(\Var(F)(1+n/N))$) equals $\min_{a,b}\Var(\hat\theta_{a+bF})=:V_{\rm PPI++}$, and
\begin{equation}\label{eq:gap}
V_{\rm PPI++}-V_{\rm eff}=\frac{N}{n(n+N)}\,\Var\big(\mu(Z)-b^\star F\big),\qquad b^\star=\Cov(Y,F)/\Var(F).
\end{equation}
\end{proposition}
\begin{corollary}[headroom and its absence]\label{cor:head}
Let $\kappa=\Var(\mu(Z)-b^\star F)/\Var(Y-b^\star F)\in[0,1]$ (the share of the residual variance of the best affine calibration that $Z$ can still explain). Then $V_{\rm eff}/V_{\rm PPI++}\to1-\kappa$ as $N/n\to\infty$, so the interval-width ratio is $\ge\sqrt{1-\kappa}$ asymptotically. (i) If $\mu(Z)$ is affine in $F$ (calibrated predictor, no informative covariates), $\kappa=0$ and PPI++ is already efficient. (ii) For finite $N/n$ the gap carries the factor $N/(n+N)$; at $N=0$ all methods reduce to the classical interval. (iii) If $\lambda^\star\notin[0,1]$, the clipped version used in \texttt{ppi\_py} is strictly worse than $V_{\rm PPI++}$ and $\kappa$ understates the headroom over it.
\end{corollary}
Proofs are in Appendix~\ref{app:proofs}. Under the model above (nonparametric $P$, labels missing completely at random) $V_{\rm eff}$ is the semiparametric efficiency bound \citep{robins1994estimation,tsiatis2006semiparametric}, hence \emph{no regular asymptotically linear estimator, and so no Wald interval built from one, can have smaller asymptotic variance}; the experiments in Section~\ref{sec:synth} confirm the formula and not just the bound. We do not claim a non-asymptotic lower bound on the width of arbitrary valid intervals.

\textbf{Feasible estimator.} Since $\mu$ is unknown we fit $\hat g$ on the labeled data with $K$-fold cross-fitting: labeled point $i$ in fold $k$ uses $\hat g^{(-k)}$ (trained only on labeled points outside fold $k$) in both terms of \eqref{eq:aug}; unlabeled points use the average of the $K$ fold models. The variance is estimated from the two samples separately,
$\hat V=n\,\widehat{\Var}_i\big[\tfrac{Y_i-\hat g_i}{n}+\tfrac{\hat g_i}{n+N}\big]+\tfrac{N}{(n+N)^2}\widehat{\Var}_j[\hat g(Z'_j)]$, and the interval is $\hat\theta\pm z_{1-\alpha/2}\sqrt{\hat V}$.
\begin{proposition}[validity, sketch]\label{prop:valid}
If $\|\hat g^{(-k)}-g^\ast\|_{L_2(P_Z)}\to_p0$ for some fixed $g^\ast$ (not necessarily $\mu$) and $\Var(Y-g^\ast)>0$, then as $n,N\to\infty$ with $N/n$ bounded away from $0$, $(\hat\theta-\theta)/\sqrt{\hat V}\Rightarrow\mathcal N(0,1)$ and the asymptotic variance is that of Proposition~\ref{prop:var} with $g=g^\ast$.
\end{proposition}
Hence validity does not require $\hat g$ to be consistent for $\mu$; consistency only buys efficiency. All of this is asymptotic, and says nothing about the finite-sample cost of estimating $\hat g$, which we measure.

\section{Methods Compared and Protocol}\label{sec:setup}
\textbf{Methods.} (1) \emph{classical}: labeled-only CLT interval; (2) \emph{PPI} ($\lambda=1$) and (3) \emph{PPI++} (power-tuned, $\lambda$ clipped to $[0,1]$), both from \texttt{ppi-python} 0.2.3 (\texttt{ppi\_mean\_ci}, which uses standard deviations with $\mathrm{ddof}=0$; we left it unchanged as the reference implementation); (4) \emph{AIPW-lin}: cross-fitted $\hat g=a+bF$; (5) \emph{AIPW-iso}: isotonic regression of $Y$ on $F$; when covariates exist also \emph{AIPW-linX} (linear in $(F,W)$) and \emph{AIPW-gbm} (histogram gradient boosting, depth 3, 60 trees \citep{pedregosa2011sklearn}); (6) \emph{AIPW-auto}: among the above, the learner with the smallest out-of-fold squared error (all candidates share folds; selection uses the same out-of-fold predictions that enter the estimator). $K=5$. All intervals are 95\% two-sided Wald intervals unless stated; Section~\ref{sec:valid} adds studentized-bootstrap versions of classical, PPI++ and AIPW-lin.

\textbf{Populations and sampling.} Each file of the PPI reference package (data version \texttt{v0.2.3}) is treated as the population; labeled and unlabeled samples are drawn i.i.d.\ \emph{with replacement}, so A1 holds exactly and the truth $\theta=\bar Y$ over the file is exact. Datasets (rows of Table~\ref{tab:main}): galaxies, forest, alphafold (covariates: three modification indicators), gene\_expression, census\_income (covariates: age, sex), census\_healthcare, census\_education, plankton (2013 half only here), ballots (the 1027 labeled ballots). Synthetic designs S0--S7 (Appendix~\ref{app:synth}) have analytic truth. Grid: $n\in\{25,50,100,200,500,1000\}$ with $N=5000$, and $N\in\{50,200,1000,5000,20000\}$ with $n=200$; \textbf{1000 replications per cell}, seeds derived from \texttt{SeedSequence}$([20261101,\text{dataset},n,N,\text{rep}])$, all methods evaluated on the same samples. Coverage has Monte-Carlo s.e.\ of about $0.7$ points near 95\%. The width ratio is mean width over mean PPI++ width on the same replications, with a paired bootstrap interval over replications (always $\le\pm0.002$ for ratios near one at $n\ge200$ and reported in the code outputs); it reflects Monte-Carlo, not dataset, uncertainty. Headroom $\kappa$ is estimated on the full population with 5-fold cross-validated learners (affine, isotonic, linear-in-$(F,W)$, boosting); since the learner class is finite, true headroom could be somewhat larger than $\hat\kappa$.

\section{Results}
\subsection{Real data: almost no headroom for the mean}\label{sec:real}
\begin{table}[t]
\caption{Coverage (\%) and width relative to PPI++ at $n=200$, $N=5000$, 1000 replications. $b$: OLS slope of $Y$ on $F$; $\hat\kappa$: estimated headroom share (Cor.~\ref{cor:head}); bound $=\sqrt{1-\hat\kappa}$ is the asymptotic ($N/n\to\infty$) width ratio of the efficient estimator \emph{relative to unclipped-$\lambda$ PPI++}. Ratio $<1$ means shorter than PPI++. Rows S$\cdot$ are synthetic. Cov.\ s.e.\ $\approx0.7$.}\label{tab:main}
\centering\scriptsize\setlength{\tabcolsep}{2.6pt}
\begin{tabular}{lrrr|ccccc|cccc}
\toprule
 & & & & \multicolumn{5}{c|}{coverage (\%)} & \multicolumn{4}{c}{width / PPI++} \\
dataset & $b$ & $\hat\kappa$ & bound & clas. & PPI & PPI++ & lin & auto & clas. & PPI & lin & auto \\
\midrule
galaxies & 0.98 & 0.000 & 1.000 & 94.1 & 93.0 & 92.5 & 93.4 & 92.8 & 1.456 & 1.001 & 1.008 & 1.006 \\
forest & 0.90 & 0.000 & 1.000 & 95.4 & 95.9 & 95.2 & 95.8 & 95.5 & 1.176 & 1.004 & 1.012 & 1.005 \\
alphafold & 1.03 & 0.008 & 0.996 & 94.8 & 95.4 & 95.6 & 95.3 & 94.2 & 1.657 & 0.995 & 1.005 & 0.998 \\
gene\_expression & 1.15 & 0.024 & 0.988 & 94.5 & 94.9 & 94.6 & 95.0 & 94.7 & 2.935 & 0.998 & 0.967 & 0.967 \\
census\_income & 1.38 & 0.011 & 0.994 & 92.4 & 90.7 & 90.6 & 90.7 & 90.6 & 1.237 & 1.000 & 1.002 & 0.995 \\
census\_healthcare & 0.98 & 0.001 & 0.999 & 96.0 & 95.4 & 95.4 & 95.3 & 95.2 & 1.390 & 1.001 & 1.010 & 1.007 \\
census\_education & 0.97 & 0.001 & 1.000 & 93.6 & 94.6 & 94.2 & 94.4 & 94.3 & 1.266 & 1.003 & 1.014 & 1.011 \\
plankton & 0.84 & 0.000 & 1.000 & 95.5 & 95.0 & 94.2 & 95.5 & 95.5 & 1.739 & 1.039 & 1.004 & 1.004 \\
ballots & 0.98 & 0.000 & 1.000 & 94.9 & 94.1 & 93.9 & 93.9 & 93.9 & 3.745 & 1.019 & 1.000 & 1.000 
\\
\midrule
S0\_calibrated & 1.00 & 0.000 & 1.000 & 94.6 & 93.5 & 93.3 & 93.7 & 93.7 & 1.387 & 1.000 & 1.010 & 1.010 \\
S1\_nonlinear\_link & 1.66 & 0.884 & 0.341 & 94.3 & 92.7 & 92.7 & 92.4 & 90.4 & 1.396 & 1.000 & 0.953 & 0.550 \\
S2\_omitted\_linear & 1.00 & 0.798 & 0.449 & 94.3 & 94.7 & 94.5 & 94.6 & 94.7 & 1.320 & 1.001 & 1.011 & 0.509 \\
S3\_omitted\_nonlinear & 1.00 & 0.806 & 0.441 & 94.0 & 95.5 & 95.3 & 95.4 & 94.9 & 1.297 & 1.001 & 1.011 & 0.605 \\
S4\_noise\_predictor & -0.00 & 0.800 & 0.447 & 95.4 & 95.4 & 95.2 & 95.4 & 96.4 & 1.001 & 1.356 & 1.012 & 0.486 \\
S5\_all\_noise & -0.01 & 0.000 & 1.000 & 94.9 & 94.4 & 95.0 & 95.1 & 94.9 & 1.001 & 1.235 & 1.012 & 1.011 \\
S6\_heavy\_tail & 1.00 & 0.000 & 1.000 & 95.9 & 95.4 & 95.2 & 95.4 & 95.3 & 1.159 & 1.001 & 1.012 & 1.011 \\
S7\_binary\_miscal & 1.49 & 0.013 & 0.993 & 94.0 & 95.2 & 95.2 & 95.3 & 94.7 & 1.228 & 1.000 & 0.983 & 0.980 
\\
\bottomrule
\end{tabular}
\end{table}
Table~\ref{tab:main} (top) shows that at $n=200$ the best cross-fitted interval is between $0.967$ and $1.011$ of the PPI++ width on all nine real datasets (median $1.004$); at $n=1000$ it lies in $[0.996,1.002]$ (Appendix Table~\ref{tab:n1000}). Estimated headroom is $\hat\kappa\le0.024$ (gene\_expression), $0.011$ (census\_income), $0.008$ (alphafold) and $\le0.002$ elsewhere, matching the near-1 ratios. Fig.~\ref{fig:ratio} (left) shows the ratio against $n$ and $N$: in nearly every real dataset the curve sits at $1.00$--$1.02$ (slightly above because estimating $\hat g$ costs a little). Increasing $N$ at $n=200$ moves the real-data ratios only for gene\_expression ($0.967$ at $N=5000$, $0.928$ at $N=20000$); all others stay within $[0.99,1.011]$.

The gene\_expression gain has a transparent source: the predictor's slope is $b=1.15>1$, so $\lambda^\star>1$ and ppi\_py clips $\lambda$ to $1$ while AIPW-lin does not. An ablation with unclipped $\lambda$ (Appendix Table~\ref{tab:noclip}) gives $0.975$ of the clipped width at $n=200$ (coverage $94.0$ vs $94.2$) and exactly $1.000$ at $n=1000$, where the finite-$N$ shrinkage $1/(1+n/N)$ brings $\lambda^\star$ below $1$ again. Unclipping alone thus recovers the gain; no calibration learner is needed. For census\_income ($b=1.38$) the same ablation gives $0.980$ at $n=200$, with coverage already low ($91.6\%$ clipped, $90.7\%$ unclipped).

\begin{figure}[t]
\centering\includegraphics[width=.82\linewidth]{../figs/fig_ratio.pdf}
\caption{Width of AIPW-auto relative to PPI++ (mean over 1000 replications; bands: paired bootstrap 95\% intervals). Real data are essentially flat at 1; synthetic designs with non-affine calibration or informative covariates (S1--S4) fall to 0.49--0.65. The ratio \emph{rises} with $n$ at fixed $N$ because the gap carries the factor $N/(n+N)$ (Eq.~\ref{eq:gap}).}\label{fig:ratio}
\end{figure}

\subsection{Synthetic designs: where headroom exists, and how it is realised}\label{sec:synth}
Table~\ref{tab:main} (bottom) and Fig.~\ref{fig:ratio} (right). S1 (monotone but non-affine calibration, $Y=e^{X_1}+$noise, $F=X_1$), S2 and S3 (the label depends on a covariate the predictor ignores, linearly/non-linearly) and S4 (a pure-noise predictor, informative covariate) have $\hat\kappa\approx0.8$--$0.88$. At $n=200$, $N=5000$ AIPW-auto is $0.550$, $0.509$, $0.605$ and $0.486$ of the PPI++ width, with coverage $90.4$, $94.7$, $94.9$ and $96.4\%$. AIPW-lin helps only on S1 ($0.953$, thanks to the unclipped slope $b=1.66$) and not on S2--S4 ($1.011$), as the theory predicts: the gain requires a calibration class that contains $\mu$, here isotonic (S1) or covariate-aware (S2--S4); the automatic selection finds it (at $n=200$ it picks isotonic in $100\%$ of replications for S1, covariate-linear for S2, boosting for S3 and covariate-linear for S4 in $99.8\%$). Where $\hat\kappa=0$ (S0 calibrated, S5 all noise, S6 heavy tails) the estimators cost $1.0$--$1.1\%$ extra width, and S7 (binary label, miscalibrated probability, $\hat\kappa=0.013$) gains $2\%$. S4 also shows PPI++'s protection: its ratio to the classical interval is $1.00$ while $\lambda=1$ PPI is $1.36$ wider than PPI++; with an informative covariate PPI++ nonetheless leaves half the width on the table.

\begin{figure}[t]
\centering
\begin{minipage}[c]{.36\linewidth}\includegraphics[width=\linewidth]{../figs/fig_pred_obs.pdf}\end{minipage}\hfill
\begin{minipage}[c]{.61\linewidth}\caption{Observed width ratio of AIPW-auto vs.\ PPI++ against the asymptotic prediction $\sqrt{V_{\rm eff}/V_{\rm PPI++}}$ computed from population quantities (Section~\ref{sec:setup}), $N=5000$, all datasets and designs, $n\in\{50,200,1000\}$. Points converge to the diagonal as $n$ grows; at small $n$ synthetic points lie above it (the cost of fitting $\hat g$ is not in the formula).}\label{fig:predobs}\end{minipage}
\end{figure}
\textbf{Theory versus experiment.} Fig.~\ref{fig:predobs} plots the observed against the predicted ratio. At $n=1000$ the synthetic points lie on the diagonal within $0.05$ (S2: predicted $0.643$, observed $0.643$; S1: $0.604$ vs.\ $0.647$). At $n=200$ the realised gain for S1 and S3 is smaller than predicted ($0.55$ vs.\ $0.40$; $0.61$ vs.\ $0.50$): flexible $\hat g$ needs data. The ratio also rises with $n$ at fixed $N$ (Fig.~\ref{fig:ratio}, second panel; S2: $0.51$ at $n=200$, $0.57$ at $n=500$, $0.64$ at $n=1000$), exactly the $N/(n+N)$ factor in \eqref{eq:gap}: labels eventually dominate the unlabeled sample and the gap closes.

\subsection{Validity at small $n$: all Wald intervals under-cover on real data}\label{sec:valid}
\begin{figure}[t]
\centering\includegraphics[width=.86\linewidth]{../figs/fig_coverage.pdf}
\caption{Empirical coverage versus labeled sample size ($N=5000$). Solid/dashed: Wald; dotted: studentized bootstrap ($B=399$, 500 replications, error bars for Wald only $\pm1.96$ s.e.). Nominal 95\% (black).}\label{fig:cov}
\end{figure}
Fig.~\ref{fig:cov} shows that Wald coverage is below nominal for \emph{every} method at small $n$ and does not reach 95\% in the heavy-tailed census\_income until $n\approx1000$. Across the nine datasets the mean coverage of PPI++ is $85.9\%$ at $n=25$, $90.0\%$ at $n=50$, $92.5\%$ at $n=100$, $94.0\%$ at $n=200$ (worst single dataset: $76.4\%$, $82.3\%$, $86.5\%$, $90.6\%$); AIPW-lin is somewhat better ($89.3$, $92.1$, $93.6$, $94.4\%$) and AIPW-auto close to PPI++ ($86.6$, $91.2$, $92.9$, $94.1\%$). In ballots (predictor error rate $0.9\%$) a sample of $25$--$100$ labels often contains no or few prediction errors and the variance estimate is too small; AIPW with a fitted $\hat g$ is worse there ($77.4\%$ at $n=25$ vs.\ $84.0\%$). A comparison of width without coverage would hence flatter whichever method under-covers: at $n=25$ AIPW-auto is $0.917$ of PPI++ on ballots but covers only $77.4\%$.

Studentized bootstrap intervals (Table~\ref{tab:boot}) raise mean coverage of PPI++ to $93.3\%$ at $n=25$ and $93.9\%$ at $n=200$ (Wald: $87.1$, $93.5$; the 500-replication s.e.\ per dataset is $\approx1$ point), while the bootstrap AIPW-lin interval is at $90.9\%$ / $93.5\%$ and is in ballots and plankton considerably worse ($81\%$ and $90\%$ at $n=25$). At $n=200$ the bootstrap AIPW-lin interval is median $1.9\%$ narrower than bootstrap PPI++ (per-dataset range $0.95$--$0.99$), with $0.4$ points lower mean coverage; at $n=25$ it is $13\%$ wider. We therefore see no robust, coverage-adjusted shortening at small $n$ on real data. The classical bootstrap interval is the best-covered ($96.0\%$ at $n=25$) but $1.1$--$1.4$ times wider than PPI++ (median over datasets).

\begin{table}[t]
\caption{Studentized-bootstrap ($B=399$, 500 replications, $N=5000$) coverage (\%) and width relative to bootstrap PPI++ (ratio of mean widths), $n=50$ (top) and $n=200$ (bottom). Coverage s.e.\ $\approx1$ point.}\label{tab:boot}
\centering\small\setlength{\tabcolsep}{4pt}
\begin{tabular}{l|ccc|cc}\toprule dataset & clas. & PPI++ & AIPW-lin & clas. width & AIPW-lin width \\\midrule
galaxies & 96.0 & 95.0 & 94.4 & 1.382 & 0.963 \\
forest & 98.0 & 95.6 & 95.6 & 0.935 & 1.126 \\
alphafold & 97.6 & 94.6 & 90.2 & 0.843 & 0.947 \\
gene\_expression & 95.0 & 95.2 & 93.0 & 2.972 & 0.848 \\
census\_income & 92.8 & 92.4 & 89.6 & 1.093 & 1.019 \\
census\_healthcare & 95.6 & 93.4 & 92.4 & 1.358 & 0.989 \\
census\_education & 94.6 & 94.4 & 93.4 & 1.209 & 0.999 \\
plankton & 97.8 & 93.6 & 91.2 & 0.980 & 1.614 \\
ballots & 96.6 & 85.6 & 80.6 & 3.311 & 0.891 
\\\midrule
galaxies & 94.6 & 94.0 & 93.0 & 1.423 & 0.981 \\
forest & 94.2 & 94.8 & 95.4 & 1.150 & 0.991 \\
alphafold & 95.0 & 93.4 & 93.2 & 1.584 & 0.956 \\
gene\_expression & 94.4 & 94.8 & 94.4 & 2.887 & 0.950 \\
census\_income & 93.2 & 91.8 & 91.0 & 1.168 & 0.987 \\
census\_healthcare & 95.2 & 95.8 & 96.6 & 1.369 & 0.993 \\
census\_education & 93.2 & 93.2 & 93.2 & 1.247 & 0.987 \\
plankton & 94.8 & 93.4 & 92.2 & 1.672 & 0.954 \\
ballots & 97.0 & 93.8 & 92.2 & 3.475 & 0.971 
\\\bottomrule\end{tabular}
\end{table}

\subsection{Failure under labeled/unlabeled shift}\label{sec:shift}
\begin{table}[t]
\caption{Plankton: labeled data from 2013, unlabeled from 2014, target the 2014 mean label ($0.071$; 2013 mean $0.134$). Coverage (\%), $N=5000$, 1000 replications.}\label{tab:shift}
\centering\small\begin{tabular}{r|ccccc}\toprule $n$ & classical & PPI & PPI++ & AIPW-lin & AIPW-auto\\\midrule
100 & 53.0 & 95.4 & 95.2 & 96.7 & 96.7 \\
500 & 1.0 & 94.7 & 94.7 & 72.1 & 72.1 \\
5000 & 0.0 & 93.0 & 0.0 & 0.0 & 0.0 

\\\bottomrule\end{tabular}
\end{table}
A1 (exchangeability) can fail: with labels from one year and unlabeled data from the next, Table~\ref{tab:shift} shows that the labeled-only interval fails (bias $0.062$), PPI with $\lambda=1$ stays valid ($93$--$95\%$; the rectifier $Y-F$ transfers across years even though $Y$ does not), while PPI++ and all cross-fitted estimators collapse to $0.0\%$ coverage at $n=5000$. Any method that uses the labeled data to learn calibration (power tuning, $\hat g$) inherits the labeled distribution's bias with a precision that grows with $n$. This is a conditional benefit: efficiency gains from fitted calibration are traded against robustness to shift, and both PPI++ and our estimators are on the fragile side.

\section{Discussion and Limitations}
\textbf{Answer to the question.} For a population mean, with a predictor that is calibrated up to an affine map and no covariate information the predictor missed, the optimally tuned PPI++ interval is already (asymptotically) as short as any regular estimator permits, and on nine real datasets we measure at most a $3.3\%$ ($n=200$, gene\_expression) reduction, explained by an un-clipped $\lambda$. The interval can be substantially shortened (here by up to $\approx50\%$) only when $\E[Y\mid Z]$ departs from an affine function of the prediction, e.g.\ non-linear calibration or unused covariates; the gain scales with $N/(n+N)$, needs enough labels for $\hat g$ (our selected learners pay off from $n\approx100$--$200$), and disappears with a calibrated predictor. Practitioners can estimate $\hat\kappa$ from the labeled data (cross-validated residual variances) before choosing a method.

\textbf{Limitations.} (i) \emph{Target.} Only the mean. Regression coefficients, quantiles and other M-estimands have analogous efficient influence functions (conditional expectation of the score given $Z$) but we did not test them; for them headroom may differ. (ii) \emph{Learner class.} $\hat\kappa$ comes from four learner families under cross-validation; richer learners might find more headroom in the real datasets, so ``at most $2.4\%$'' is a statement about these learners. (iii) \emph{Finite samples.} Proposition~\ref{prop:valid} is asymptotic and its sketch not a full proof; at $n\le100$ real-data coverage falls below nominal for all methods and our bootstrap fix is partial. We did not try other corrections (Edgeworth, exact binomial); $t$-quantile versions are stored in the result files but not analysed. (iv) \emph{Sampling design.} Datasets are resampled with replacement from a finite file, so results are about the empirical distribution of each file; small files (forest, ballots) are heavily bootstrapped at $N=5000$. The ballots population is the 1027 labeled ballots only. (v) \emph{Exchangeability} is assumed; Section~\ref{sec:shift} is a single shift example. (vi) The classical and PPI baselines use the \texttt{ppi-python} standard errors ($\mathrm{ddof}=0$), our estimators $\mathrm{ddof}=1$, a $2\%$ width difference at $n=25$ that we did not remove. (vii) We did not study quantile/binary-optimal intervals such as Wilson or exact methods, which may be shorter at small $n$ for rare-error binary labels. (viii) Compute was small (CPU only) and reproducibility is in the supplementary code.

\textbf{Broader impact.} The work is methodological and uses public data; the main practical risk is overconfidence from intervals reported in the small-label regime where we find under-coverage.

\bibliographystyle{plainnat}
\bibliography{refs}

\appendix
\section{Proofs}\label{app:proofs}
\textbf{Proposition~\ref{prop:var}.} Write $\hat\theta_g=\sum_{i\le n}\big[\tfrac{Y_i}{n}+(\tfrac1{n+N}-\tfrac1n)g(Z_i)\big]+\tfrac1{n+N}\sum_{j\le N}g(Z_j')$; the two sums are independent sums of i.i.d.\ terms, $c=\tfrac1{n+N}$, $a=\tfrac1n$. Then $\Var=n[a^2\Var Y+2a(c-a)\Cov(Y,g)+(c-a)^2\Var g]+Nc^2\Var g$. With $c-a=-\tfrac N{n(n+N)}$: $n(c-a)^2+Nc^2=\tfrac N{n(n+N)}$ and $2na(c-a)=-\tfrac{2N}{n(n+N)}$, so $\Var=\tfrac{\Var Y}n+\tfrac N{n(n+N)}(\Var g-2\Cov(Y,g))=\tfrac{\Var Y}{n+N}+\tfrac N{n(n+N)}\Var(Y-g)$, using $\tfrac1n-\tfrac N{n(n+N)}=\tfrac1{n+N}$. Unbiasedness: $\E g$ cancels between the terms. The minimum over $g$ of $\Var(Y-g(Z))$ is $\E\Var(Y\mid Z)$ at $g=\mu$.\\
\textbf{Proposition~\ref{prop:ppi}.} $\Var(Y-\lambda F)/n+\lambda^2\Var F/N$ is a quadratic in $\lambda$ minimized at $\lambda^\star$ with value $\frac1n[\Var Y-\Cov(Y,F)^2/(\Var F(1+n/N))]=\frac1n[\Var Y-\tfrac{N}{n+N}b^{\star2}\Var F]$. By Prop.~\ref{prop:var}, $\min_b\Var(\hat\theta_{bF})=\frac{\Var Y}{n+N}+\frac N{n(n+N)}(\Var Y-b^{\star2}\Var F)$, which equals the same expression. For \eqref{eq:gap}: $\Var(Y-b^\star F)=\E\Var(Y\mid Z)+\Var(\mu-b^\star F)$ because $Y-\mu\perp(\mu-b^\star F)$; subtract $V_{\rm eff}$.\\
\textbf{Corollary~\ref{cor:head}.} As $N/n\to\infty$, $nV_{\rm PPI++}\to\Var(Y-b^\star F)$ and $nV_{\rm eff}\to\E\Var(Y\mid Z)$, whose ratio is $1-\kappa$. (i) $\mu$ affine in $F$ implies $\mu=a+b^\star F$ and the gap vanishes. (iii) $\lambda\mapsto\Var(\hat\theta^{\rm PPI}_\lambda)$ is strictly convex, so clipping an interior minimizer is strictly suboptimal.\\
\textbf{Proposition~\ref{prop:valid} (sketch).} Let $\hat\theta^\ast$ denote \eqref{eq:aug} with $g=g^\ast$. For fold $k$, conditional on the training folds, $\hat\theta-\hat\theta^\ast$ is a sum over folds of $(\text{mean}_{\rm all}-\text{mean}_{\rm lab})(\hat g^{(-k)}-g^\ast)$, a centred difference of sample means of one fixed function with variance $\le C\,\|\hat g^{(-k)}-g^\ast\|^2_{L_2}\,(1/n)=o_p(1/n)$ (the usual double machine learning argument \citep{chernozhukov2018dml}); the unlabeled average over the $K$ fold models adds only terms of the same type. Then $\hat\theta^\ast$ is asymptotically normal by the Lindeberg CLT applied to the two independent samples with variance from Prop.~\ref{prop:var}, and $\hat V$ is consistent because the empirical variances of the plug-in terms converge to those of the $g^\ast$-terms by the same $L_2$ argument. A full proof with explicit moment conditions is not given.

\section{Synthetic designs}\label{app:synth}
$X_1,X_2,F_0\sim\mathcal N(0,1)$, $\varepsilon\sim\mathcal N(0,1)$ independent; $W$ are the covariates available to the analyst. \textbf{S0:} $Y=X_1+\varepsilon$, $F=X_1$. \textbf{S1:} $Y=e^{X_1}+0.5\varepsilon$, $F=X_1$ ($\theta=e^{1/2}$). \textbf{S2:} $Y=X_1+X_2+0.5\varepsilon$, $F=X_1$, $W=X_2$. \textbf{S3:} $Y=X_1+1.5\sin(2X_2)+0.5\varepsilon$, $F=X_1$, $W=X_2$. \textbf{S4:} $Y=X_1+0.5\varepsilon$, $F=F_0$ (noise), $W=X_1$. \textbf{S5:} $Y=X_1+\varepsilon$, $F=F_0$, $W$ two independent noise columns. \textbf{S6:} $Y=X_1+t_3$, $F=X_1$. \textbf{S7:} $Y\sim\mathrm{Bernoulli}(\sigma(2X_1))$, $F=\sigma(X_1)$ ($\theta=0.5$). Population constants ($\hat\kappa$, $b$) are from $2\times10^5$ draws.

\section{Additional results}
\begin{table}[h]
\caption{As Table~\ref{tab:main} for $n=1000$ (top: real; bottom: synthetic).}\label{tab:n1000}
\centering\scriptsize\setlength{\tabcolsep}{2.6pt}
\begin{tabular}{lrrr|ccccc|cccc}\toprule
 & & & & \multicolumn{5}{c|}{coverage (\%)} & \multicolumn{4}{c}{width / PPI++} \\
dataset & $b$ & $\hat\kappa$ & bound & clas. & PPI & PPI++ & lin & auto & clas. & PPI & lin & auto \\\midrule
galaxies & 0.98 & 0.000 & 1.000 & 94.7 & 93.9 & 93.4 & 93.5 & 93.4 & 1.358 & 1.021 & 1.001 & 1.001 \\
forest & 0.90 & 0.000 & 1.000 & 95.5 & 94.6 & 94.9 & 95.0 & 94.4 & 1.148 & 1.018 & 1.002 & 0.996 \\
alphafold & 1.03 & 0.008 & 0.996 & 94.2 & 96.0 & 95.9 & 95.8 & 95.2 & 1.498 & 1.015 & 1.001 & 0.999 \\
gene\_expression & 1.15 & 0.024 & 0.988 & 95.2 & 95.3 & 94.8 & 95.1 & 95.1 & 2.099 & 1.002 & 1.000 & 1.000 \\
census\_income & 1.38 & 0.011 & 0.994 & 95.8 & 95.0 & 95.0 & 94.6 & 94.5 & 1.199 & 1.000 & 1.000 & 0.998 \\
census\_healthcare & 0.98 & 0.001 & 0.999 & 94.8 & 94.6 & 94.7 & 95.1 & 95.2 & 1.308 & 1.017 & 1.002 & 1.001 \\
census\_education & 0.97 & 0.001 & 1.000 & 95.7 & 96.3 & 95.7 & 96.0 & 96.2 & 1.220 & 1.013 & 1.002 & 1.002 \\
plankton & 0.84 & 0.000 & 1.000 & 95.2 & 95.3 & 94.9 & 95.5 & 95.5 & 1.548 & 1.116 & 1.001 & 1.001 \\
ballots & 0.98 & 0.000 & 1.000 & 95.1 & 94.9 & 95.4 & 95.4 & 95.4 & 2.253 & 1.101 & 1.000 & 1.000 
\\\midrule
S0\_calibrated & 1.00 & 0.000 & 1.000 & 95.0 & 95.4 & 95.9 & 95.8 & 95.8 & 1.309 & 1.014 & 1.002 & 1.002 \\
S1\_nonlinear\_link & 1.66 & 0.884 & 0.341 & 94.3 & 94.4 & 94.4 & 94.6 & 93.6 & 1.330 & 1.000 & 0.974 & 0.647 \\
S2\_omitted\_linear & 1.00 & 0.798 & 0.449 & 94.6 & 94.3 & 94.7 & 94.7 & 94.6 & 1.260 & 1.012 & 1.002 & 0.643 \\
S3\_omitted\_nonlinear & 1.00 & 0.806 & 0.441 & 95.1 & 95.5 & 95.0 & 95.0 & 94.5 & 1.241 & 1.011 & 1.002 & 0.650 \\
S4\_noise\_predictor & -0.00 & 0.800 & 0.447 & 95.1 & 95.5 & 95.2 & 95.0 & 95.2 & 1.000 & 1.402 & 1.002 & 0.579 \\
S5\_all\_noise & -0.01 & 0.000 & 1.000 & 94.3 & 95.3 & 94.3 & 94.3 & 94.4 & 1.000 & 1.264 & 1.002 & 1.002 \\
S6\_heavy\_tail & 1.00 & 0.000 & 1.000 & 95.4 & 95.6 & 95.4 & 95.8 & 95.7 & 1.129 & 1.006 & 1.002 & 1.002 \\
S7\_binary\_miscal & 1.49 & 0.013 & 0.993 & 94.2 & 94.5 & 94.5 & 93.8 & 94.0 & 1.203 & 1.000 & 0.993 & 0.991 
\\\bottomrule\end{tabular}
\end{table}
\begin{table}[h]
\caption{As Table~\ref{tab:main} for $n=50$.}\label{tab:n50}
\centering\scriptsize\setlength{\tabcolsep}{2.6pt}
\begin{tabular}{lrrr|ccccc|cccc}\toprule
 & & & & \multicolumn{5}{c|}{coverage (\%)} & \multicolumn{4}{c}{width / PPI++} \\
dataset & $b$ & $\hat\kappa$ & bound & clas. & PPI & PPI++ & lin & auto & clas. & PPI & lin & auto \\\midrule
galaxies & 0.98 & 0.000 & 1.000 & 93.9 & 95.3 & 93.5 & 95.4 & 94.4 & 1.486 & 0.996 & 1.040 & 1.025 \\
forest & 0.90 & 0.000 & 1.000 & 94.0 & 91.9 & 90.4 & 93.5 & 92.2 & 1.192 & 1.007 & 1.064 & 1.048 \\
alphafold & 1.03 & 0.008 & 0.996 & 94.6 & 92.7 & 88.9 & 93.6 & 90.2 & 1.723 & 0.982 & 1.030 & 0.991 \\
gene\_expression & 1.15 & 0.024 & 0.988 & 93.6 & 94.4 & 93.7 & 96.0 & 95.7 & 3.220 & 0.976 & 0.933 & 0.929 \\
census\_income & 1.38 & 0.011 & 0.994 & 88.0 & 83.4 & 82.3 & 84.0 & 82.9 & 1.270 & 0.998 & 1.055 & 1.032 \\
census\_healthcare & 0.98 & 0.001 & 0.999 & 96.1 & 94.0 & 93.3 & 94.9 & 93.8 & 1.413 & 1.000 & 1.044 & 1.030 \\
census\_education & 0.97 & 0.001 & 1.000 & 92.6 & 91.8 & 90.4 & 92.4 & 92.0 & 1.291 & 1.006 & 1.056 & 1.038 \\
plankton & 0.84 & 0.000 & 1.000 & 90.4 & 96.8 & 90.5 & 94.9 & 94.9 & 1.893 & 1.005 & 1.051 & 1.051 \\
ballots & 0.98 & 0.000 & 1.000 & 94.4 & 84.3 & 87.3 & 84.5 & 84.5 & 5.397 & 0.950 & 0.971 & 0.971 
\\\midrule
S0\_calibrated & 1.00 & 0.000 & 1.000 & 93.9 & 94.4 & 94.2 & 95.2 & 95.1 & 1.409 & 0.998 & 1.044 & 1.039 \\
S1\_nonlinear\_link & 1.66 & 0.884 & 0.341 & 89.0 & 84.9 & 84.5 & 88.4 & 86.4 & 1.445 & 0.997 & 1.005 & 0.674 \\
S2\_omitted\_linear & 1.00 & 0.798 & 0.449 & 94.2 & 93.9 & 94.0 & 94.7 & 95.0 & 1.340 & 0.999 & 1.048 & 0.491 \\
S3\_omitted\_nonlinear & 1.00 & 0.806 & 0.441 & 93.0 & 93.9 & 93.3 & 94.0 & 93.2 & 1.309 & 1.001 & 1.046 & 0.894 \\
S4\_noise\_predictor & -0.00 & 0.800 & 0.447 & 93.4 & 93.6 & 93.3 & 94.2 & 94.6 & 1.005 & 1.359 & 1.051 & 0.488 \\
S5\_all\_noise & -0.01 & 0.000 & 1.000 & 95.8 & 95.0 & 95.1 & 95.9 & 95.0 & 1.005 & 1.232 & 1.051 & 1.041 \\
S6\_heavy\_tail & 1.00 & 0.000 & 1.000 & 95.2 & 95.1 & 94.8 & 95.4 & 95.0 & 1.183 & 1.002 & 1.048 & 1.042 \\
S7\_binary\_miscal & 1.49 & 0.013 & 0.993 & 93.1 & 93.5 & 93.5 & 93.9 & 92.4 & 1.234 & 1.000 & 1.006 & 0.997 
\\\bottomrule\end{tabular}
\end{table}
\begin{table}[h]
\caption{PPI++ with $\lambda$ clipped to $[0,1]$ (as in \texttt{ppi\_py}) versus unclipped; $N=5000$, 2000 replications. Columns: dataset, $n$, coverage clipped, coverage unclipped (\%), width ratio unclipped/clipped.}\label{tab:noclip}
\centering\small\begin{tabular}{lrrrr}\toprule dataset & $n$ & cov.\ clipped & cov.\ unclipped & width ratio \\\midrule
S1\_nonlinear\_link & 200 & 91.8 & 91.5 & 0.924 \\
S1\_nonlinear\_link & 1000 & 92.9 & 93.0 & 0.969 \\
S7\_binary\_miscal & 200 & 94.5 & 94.5 & 0.974 \\
S7\_binary\_miscal & 1000 & 94.4 & 94.7 & 0.991 \\
alphafold & 200 & 93.8 & 93.4 & 1.008 \\
alphafold & 1000 & 94.6 & 94.6 & 1.000 \\
census\_income & 200 & 91.6 & 90.7 & 0.980 \\
census\_income & 1000 & 94.5 & 94.1 & 0.996 \\
gene\_expression & 200 & 94.2 & 94.0 & 0.975 \\
gene\_expression & 1000 & 95.4 & 95.4 & 1.000 
\\\bottomrule\end{tabular}
\end{table}

\section{Reproducibility and notes}
Code: \texttt{code/} (run order and commands in \texttt{README.md}); data: \texttt{ppi-python} 0.2.3 example sets (downloaded via \texttt{ppi\_py.datasets.load\_dataset}); Python 3.11, numpy 2.4.6, scipy 1.17.1, pandas 3.1.0, scikit-learn 1.9.1, statsmodels 0.15.0 (installed, unused), matplotlib 3.11.2. Main runs took about 42 min (real), 41 min (synthetic), 24 min (bootstrap) wall-clock on a shared 40-core machine (observed per-process CPU share suggests an effective throughput of roughly 4--5 cores). No GPU, no paid API. No part of this paper's analysis depends on the plankton 2014 labels except Section~\ref{sec:shift}. Table~\ref{tab:main}/\ref{tab:n1000}/\ref{tab:n50}: \texttt{run\_experiments.py}, \texttt{analysis.py}, \texttt{theory.py}, \texttt{make\_report.py}; Fig.~\ref{fig:ratio}, \ref{fig:predobs}: same; Fig.~\ref{fig:cov}, Table~\ref{tab:boot}: \texttt{run\_boot.py}; Table~\ref{tab:shift}: \texttt{run\_shift.py}; Table~\ref{tab:noclip}: \texttt{run\_noclip.py}.

\newpage
\section*{NeurIPS Paper Checklist}




\begin{enumerate}

\item {\bf Claims}
    \item[] Question: Do the main claims made in the abstract and introduction accurately reflect the paper's contributions and scope?
    \item[] Answer: \answerYes{} % Replace by \answerYes{}, \answerNo{}, or \answerNA{}.
    \item[] Justification: Abstract and Section 1 state a conditional answer (little headroom for the mean on real data, large headroom in designs with non-affine calibration or unused covariates, validity caveats), each backed by Sections 3--5.
    \item[] Guidelines:
    \begin{itemize}
        \item The answer \answerNA{} means that the abstract and introduction do not include the claims made in the paper.
        \item The abstract and/or introduction should clearly state the claims made, including the contributions made in the paper and important assumptions and limitations. A \answerNo{} or \answerNA{} answer to this question will not be perceived well by the reviewers. 
        \item The claims made should match theoretical and experimental results, and reflect how much the results can be expected to generalize to other settings. 
        \item It is fine to include aspirational goals as motivation as long as it is clear that these goals are not attained by the paper. 
    \end{itemize}

\item {\bf Limitations}
    \item[] Question: Does the paper discuss the limitations of the work performed by the authors?
    \item[] Answer: \answerYes{} % Replace by \answerYes{}, \answerNo{}, or \answerNA{}.
    \item[] Justification: Section 6 lists limitations: target, learner class, asymptotic theory, finite-sample under-coverage, resampling design, exchangeability, baseline variance conventions.
    \item[] Guidelines:
    \begin{itemize}
        \item The answer \answerNA{} means that the paper has no limitation while the answer \answerNo{} means that the paper has limitations, but those are not discussed in the paper. 
        \item The authors are encouraged to create a separate ``Limitations'' section in their paper.
        \item The paper should point out any strong assumptions and how robust the results are to violations of these assumptions (e.g., independence assumptions, noiseless settings, model well-specification, asymptotic approximations only holding locally). The authors should reflect on how these assumptions might be violated in practice and what the implications would be.
        \item The authors should reflect on the scope of the claims made, e.g., if the approach was only tested on a few datasets or with a few runs. In general, empirical results often depend on implicit assumptions, which should be articulated.
        \item The authors should reflect on the factors that influence the performance of the approach. For example, a facial recognition algorithm may perform poorly when image resolution is low or images are taken in low lighting. Or a speech-to-text system might not be used reliably to provide closed captions for online lectures because it fails to handle technical jargon.
        \item The authors should discuss the computational efficiency of the proposed algorithms and how they scale with dataset size.
        \item If applicable, the authors should discuss possible limitations of their approach to address problems of privacy and fairness.
        \item While the authors might fear that complete honesty about limitations might be used by reviewers as grounds for rejection, a worse outcome might be that reviewers discover limitations that aren't acknowledged in the paper. The authors should use their best judgment and recognize that individual actions in favor of transparency play an important role in developing norms that preserve the integrity of the community. Reviewers will be specifically instructed to not penalize honesty concerning limitations.
    \end{itemize}

\item {\bf Theory assumptions and proofs}
    \item[] Question: For each theoretical result, does the paper provide the full set of assumptions and a complete (and correct) proof?
    \item[] Answer: \answerYes{} % Replace by \answerYes{}, \answerNo{}, or \answerNA{}.
    \item[] Justification: Section 3 states assumptions; proofs of Props.~1--2 and Cor.~1 are in Appendix A; Prop.~3 is given only as a sketch and labeled as such.
    \item[] Guidelines:
    \begin{itemize}
        \item The answer \answerNA{} means that the paper does not include theoretical results. 
        \item All the theorems, formulas, and proofs in the paper should be numbered and cross-referenced.
        \item All assumptions should be clearly stated or referenced in the statement of any theorems.
        \item The proofs can either appear in the main paper or the supplemental material, but if they appear in the supplemental material, the authors are encouraged to provide a short proof sketch to provide intuition. 
        \item Inversely, any informal proof provided in the core of the paper should be complemented by formal proofs provided in appendix or supplemental material.
        \item Theorems and Lemmas that the proof relies upon should be properly referenced. 
    \end{itemize}

    \item {\bf Experimental result reproducibility}
    \item[] Question: Does the paper fully disclose all the information needed to reproduce the main experimental results of the paper to the extent that it affects the main claims and/or conclusions of the paper (regardless of whether the code and data are provided or not)?
    \item[] Answer: \answerYes{} % Replace by \answerYes{}, \answerNo{}, or \answerNA{}.
    \item[] Justification: Section 4 and Appendices B, D describe estimators, learners, grid, seeds, populations, synthetic designs; code reproduces every table and figure.
    \item[] Guidelines:
    \begin{itemize}
        \item The answer \answerNA{} means that the paper does not include experiments.
        \item If the paper includes experiments, a \answerNo{} answer to this question will not be perceived well by the reviewers: Making the paper reproducible is important, regardless of whether the code and data are provided or not.
        \item If the contribution is a dataset and\slash or model, the authors should describe the steps taken to make their results reproducible or verifiable. 
        \item Depending on the contribution, reproducibility can be accomplished in various ways. For example, if the contribution is a novel architecture, describing the architecture fully might suffice, or if the contribution is a specific model and empirical evaluation, it may be necessary to either make it possible for others to replicate the model with the same dataset, or provide access to the model. In general. releasing code and data is often one good way to accomplish this, but reproducibility can also be provided via detailed instructions for how to replicate the results, access to a hosted model (e.g., in the case of a large language model), releasing of a model checkpoint, or other means that are appropriate to the research performed.
        \item While NeurIPS does not require releasing code, the conference does require all submissions to provide some reasonable avenue for reproducibility, which may depend on the nature of the contribution. For example
        \begin{enumerate}
            \item If the contribution is primarily a new algorithm, the paper should make it clear how to reproduce that algorithm.
            \item If the contribution is primarily a new model architecture, the paper should describe the architecture clearly and fully.
            \item If the contribution is a new model (e.g., a large language model), then there should either be a way to access this model for reproducing the results or a way to reproduce the model (e.g., with an open-source dataset or instructions for how to construct the dataset).
            \item We recognize that reproducibility may be tricky in some cases, in which case authors are welcome to describe the particular way they provide for reproducibility. In the case of closed-source models, it may be that access to the model is limited in some way (e.g., to registered users), but it should be possible for other researchers to have some path to reproducing or verifying the results.
        \end{enumerate}
    \end{itemize}


\item {\bf Open access to data and code}
    \item[] Question: Does the paper provide open access to the data and code, with sufficient instructions to faithfully reproduce the main experimental results, as described in supplemental material?
    \item[] Answer: \answerYes{} % Replace by \answerYes{}, \answerNo{}, or \answerNA{}.
    \item[] Justification: Code and run order are supplied in the submission (code/README.md); data are the public ppi-python 0.2.3 example sets, downloaded by a documented script call.
    \item[] Guidelines:
    \begin{itemize}
        \item The answer \answerNA{} means that paper does not include experiments requiring code.
        \item Please see the NeurIPS code and data submission guidelines (\url{https://neurips.cc/public/guides/CodeSubmissionPolicy}) for more details.
        \item While we encourage the release of code and data, we understand that this might not be possible, so \answerNo{} is an acceptable answer. Papers cannot be rejected simply for not including code, unless this is central to the contribution (e.g., for a new open-source benchmark).
        \item The instructions should contain the exact command and environment needed to run to reproduce the results. See the NeurIPS code and data submission guidelines (\url{https://neurips.cc/public/guides/CodeSubmissionPolicy}) for more details.
        \item The authors should provide instructions on data access and preparation, including how to access the raw data, preprocessed data, intermediate data, and generated data, etc.
        \item The authors should provide scripts to reproduce all experimental results for the new proposed method and baselines. If only a subset of experiments are reproducible, they should state which ones are omitted from the script and why.
        \item At submission time, to preserve anonymity, the authors should release anonymized versions (if applicable).
        \item Providing as much information as possible in supplemental material (appended to the paper) is recommended, but including URLs to data and code is permitted.
    \end{itemize}


\item {\bf Experimental setting/details}
    \item[] Question: Does the paper specify all the training and test details (e.g., data splits, hyperparameters, how they were chosen, type of optimizer) necessary to understand the results?
    \item[] Answer: \answerYes{} % Replace by \answerYes{}, \answerNo{}, or \answerNA{}.
    \item[] Justification: Section 4: folds, learner hyperparameters, selection rule, grids, seeds, number of replications; no neural training is involved.
    \item[] Guidelines:
    \begin{itemize}
        \item The answer \answerNA{} means that the paper does not include experiments.
        \item The experimental setting should be presented in the core of the paper to a level of detail that is necessary to appreciate the results and make sense of them.
        \item The full details can be provided either with the code, in appendix, or as supplemental material.
    \end{itemize}

\item {\bf Experiment statistical significance}
    \item[] Question: Does the paper report error bars suitably and correctly defined or other appropriate information about the statistical significance of the experiments?
    \item[] Answer: \answerYes{} % Replace by \answerYes{}, \answerNo{}, or \answerNA{}.
    \item[] Justification: Coverage reported with Monte-Carlo standard errors, width ratios with paired-bootstrap intervals (in result files), 1000 (500 for bootstrap-t) replications per cell; Figures show error bars/bands.
    \item[] Guidelines:
    \begin{itemize}
        \item The answer \answerNA{} means that the paper does not include experiments.
        \item The authors should answer \answerYes{} if the results are accompanied by error bars, confidence intervals, or statistical significance tests, at least for the experiments that support the main claims of the paper.
        \item The factors of variability that the error bars are capturing should be clearly stated (for example, train/test split, initialization, random drawing of some parameter, or overall run with given experimental conditions).
        \item The method for calculating the error bars should be explained (closed form formula, call to a library function, bootstrap, etc.)
        \item The assumptions made should be given (e.g., Normally distributed errors).
        \item It should be clear whether the error bar is the standard deviation or the standard error of the mean.
        \item It is OK to report 1-sigma error bars, but one should state it. The authors should preferably report a 2-sigma error bar than state that they have a 96\% CI, if the hypothesis of Normality of errors is not verified.
        \item For asymmetric distributions, the authors should be careful not to show in tables or figures symmetric error bars that would yield results that are out of range (e.g., negative error rates).
        \item If error bars are reported in tables or plots, the authors should explain in the text how they were calculated and reference the corresponding figures or tables in the text.
    \end{itemize}

\item {\bf Experiments compute resources}
    \item[] Question: For each experiment, does the paper provide sufficient information on the computer resources (type of compute workers, memory, time of execution) needed to reproduce the experiments?
    \item[] Answer: \answerYes{} % Replace by \answerYes{}, \answerNo{}, or \answerNA{}.
    \item[] Justification: Appendix D: CPU-only, wall-clock times per experiment, machine description; total well under 8 CPU-hours on a shared machine.
    \item[] Guidelines:
    \begin{itemize}
        \item The answer \answerNA{} means that the paper does not include experiments.
        \item The paper should indicate the type of compute workers CPU or GPU, internal cluster, or cloud provider, including relevant memory and storage.
        \item The paper should provide the amount of compute required for each of the individual experimental runs as well as estimate the total compute. 
        \item The paper should disclose whether the full research project required more compute than the experiments reported in the paper (e.g., preliminary or failed experiments that didn't make it into the paper). 
    \end{itemize}
    
\item {\bf Code of ethics}
    \item[] Question: Does the research conducted in the paper conform, in every respect, with the NeurIPS Code of Ethics \url{https://neurips.cc/public/EthicsGuidelines}?
    \item[] Answer: \answerYes{} % Replace by \answerYes{}, \answerNo{}, or \answerNA{}.
    \item[] Justification: The work uses public aggregate datasets and simulation; no ethics concerns arise.
    \item[] Guidelines:
    \begin{itemize}
        \item The answer \answerNA{} means that the authors have not reviewed the NeurIPS Code of Ethics.
        \item If the authors answer \answerNo, they should explain the special circumstances that require a deviation from the Code of Ethics.
        \item The authors should make sure to preserve anonymity (e.g., if there is a special consideration due to laws or regulations in their jurisdiction).
    \end{itemize}


\item {\bf Broader impacts}
    \item[] Question: Does the paper discuss both potential positive societal impacts and negative societal impacts of the work performed?
    \item[] Answer: \answerYes{} % Replace by \answerYes{}, \answerNo{}, or \answerNA{}.
    \item[] Justification: Brief statement at the end of Section 6: methodological work; the risk is overconfident intervals in small-label regimes.
    \item[] Guidelines:
    \begin{itemize}
        \item The answer \answerNA{} means that there is no societal impact of the work performed.
        \item If the authors answer \answerNA{} or \answerNo, they should explain why their work has no societal impact or why the paper does not address societal impact.
        \item Examples of negative societal impacts include potential malicious or unintended uses (e.g., disinformation, generating fake profiles, surveillance), fairness considerations (e.g., deployment of technologies that could make decisions that unfairly impact specific groups), privacy considerations, and security considerations.
        \item The conference expects that many papers will be foundational research and not tied to particular applications, let alone deployments. However, if there is a direct path to any negative applications, the authors should point it out. For example, it is legitimate to point out that an improvement in the quality of generative models could be used to generate Deepfakes for disinformation. On the other hand, it is not needed to point out that a generic algorithm for optimizing neural networks could enable people to train models that generate Deepfakes faster.
        \item The authors should consider possible harms that could arise when the technology is being used as intended and functioning correctly, harms that could arise when the technology is being used as intended but gives incorrect results, and harms following from (intentional or unintentional) misuse of the technology.
        \item If there are negative societal impacts, the authors could also discuss possible mitigation strategies (e.g., gated release of models, providing defenses in addition to attacks, mechanisms for monitoring misuse, mechanisms to monitor how a system learns from feedback over time, improving the efficiency and accessibility of ML).
    \end{itemize}
    
\item {\bf Safeguards}
    \item[] Question: Does the paper describe safeguards that have been put in place for responsible release of data or models that have a high risk for misuse (e.g., pre-trained language models, image generators, or scraped datasets)?
    \item[] Answer: \answerNA{} % Replace by \answerYes{}, \answerNo{}, or \answerNA{}.
    \item[] Justification: No models or data with misuse risk are released.
    \item[] Guidelines:
    \begin{itemize}
        \item The answer \answerNA{} means that the paper poses no such risks.
        \item Released models that have a high risk for misuse or dual-use should be released with necessary safeguards to allow for controlled use of the model, for example by requiring that users adhere to usage guidelines or restrictions to access the model or implementing safety filters. 
        \item Datasets that have been scraped from the Internet could pose safety risks. The authors should describe how they avoided releasing unsafe images.
        \item We recognize that providing effective safeguards is challenging, and many papers do not require this, but we encourage authors to take this into account and make a best faith effort.
    \end{itemize}

\item {\bf Licenses for existing assets}
    \item[] Question: Are the creators or original owners of assets (e.g., code, data, models), used in the paper, properly credited and are the license and terms of use explicitly mentioned and properly respected?
    \item[] Answer: \answerYes{} % Replace by \answerYes{}, \answerNo{}, or \answerNA{}.
    \item[] Justification: Data: ppi-python example sets (MIT-licensed code; the repository states no separate license for the .npz files, noted in README); sklearn, scipy, numpy, pandas, matplotlib cited/credited.
    \item[] Guidelines:
    \begin{itemize}
        \item The answer \answerNA{} means that the paper does not use existing assets.
        \item The authors should cite the original paper that produced the code package or dataset.
        \item The authors should state which version of the asset is used and, if possible, include a URL.
        \item The name of the license (e.g., CC-BY 4.0) should be included for each asset.
        \item For scraped data from a particular source (e.g., website), the copyright and terms of service of that source should be provided.
        \item If assets are released, the license, copyright information, and terms of use in the package should be provided. For popular datasets, \url{paperswithcode.com/datasets} has curated licenses for some datasets. Their licensing guide can help determine the license of a dataset.
        \item For existing datasets that are re-packaged, both the original license and the license of the derived asset (if it has changed) should be provided.
        \item If this information is not available online, the authors are encouraged to reach out to the asset's creators.
    \end{itemize}

\item {\bf New assets}
    \item[] Question: Are new assets introduced in the paper well documented and is the documentation provided alongside the assets?
    \item[] Answer: \answerYes{} % Replace by \answerYes{}, \answerNo{}, or \answerNA{}.
    \item[] Justification: The code released with the paper is documented in code/README.md.
    \item[] Guidelines:
    \begin{itemize}
        \item The answer \answerNA{} means that the paper does not release new assets.
        \item Researchers should communicate the details of the dataset\slash code\slash model as part of their submissions via structured templates. This includes details about training, license, limitations, etc. 
        \item The paper should discuss whether and how consent was obtained from people whose asset is used.
        \item At submission time, remember to anonymize your assets (if applicable). You can either create an anonymized URL or include an anonymized zip file.
    \end{itemize}

\item {\bf Crowdsourcing and research with human subjects}
    \item[] Question: For crowdsourcing experiments and research with human subjects, does the paper include the full text of instructions given to participants and screenshots, if applicable, as well as details about compensation (if any)? 
    \item[] Answer: \answerNA{} % Replace by \answerYes{}, \answerNo{}, or \answerNA{}.
    \item[] Justification: No crowdsourcing or human-subject research was conducted by the authors (datasets contain pre-existing human labels).
    \item[] Guidelines:
    \begin{itemize}
        \item The answer \answerNA{} means that the paper does not involve crowdsourcing nor research with human subjects.
        \item Including this information in the supplemental material is fine, but if the main contribution of the paper involves human subjects, then as much detail as possible should be included in the main paper. 
        \item According to the NeurIPS Code of Ethics, workers involved in data collection, curation, or other labor should be paid at least the minimum wage in the country of the data collector. 
    \end{itemize}

\item {\bf Institutional review board (IRB) approvals or equivalent for research with human subjects}
    \item[] Question: Does the paper describe potential risks incurred by study participants, whether such risks were disclosed to the subjects, and whether Institutional Review Board (IRB) approvals (or an equivalent approval/review based on the requirements of your country or institution) were obtained?
    \item[] Answer: \answerNA{} % Replace by \answerYes{}, \answerNo{}, or \answerNA{}.
    \item[] Justification: No human-subject research was conducted.
    \item[] Guidelines:
    \begin{itemize}
        \item The answer \answerNA{} means that the paper does not involve crowdsourcing nor research with human subjects.
        \item Depending on the country in which research is conducted, IRB approval (or equivalent) may be required for any human subjects research. If you obtained IRB approval, you should clearly state this in the paper. 
        \item We recognize that the procedures for this may vary significantly between institutions and locations, and we expect authors to adhere to the NeurIPS Code of Ethics and the guidelines for their institution. 
        \item For initial submissions, do not include any information that would break anonymity (if applicable), such as the institution conducting the review.
    \end{itemize}

\item {\bf Declaration of LLM usage}
    \item[] Question: Does the paper describe the usage of LLMs if it is an important, original, or non-standard component of the core methods in this research? Note that if the LLM is used only for writing, editing, or formatting purposes and does \emph{not} impact the core methodology, scientific rigor, or originality of the research, declaration is not required.
    %this research? 
    \item[] Answer: \answerNA{} % Replace by \answerYes{}, \answerNo{}, or \answerNA{}.
    \item[] Justification: LLMs are not part of the method; a language-model agent assisted with coding and drafting, which is not a core methodological component.
    \item[] Guidelines:
    \begin{itemize}
        \item The answer \answerNA{} means that the core method development in this research does not involve LLMs as any important, original, or non-standard components.
        \item Please refer to our LLM policy in the NeurIPS handbook for what should or should not be described.
    \end{itemize}

\end{enumerate}
\end{document}
